% !TeX spellcheck = es_PE
% document configuration
% document class
\documentclass[10pt]{beamer}

% presentation theme
\usetheme{Madrid}

% libraries
\usepackage[spanish]{babel}
\usepackage[utf8]{inputenc} % special symbols
\usepackage{xcolor} % color
\usepackage{multicol} % multiple columns
\usepackage{multirow} % multiple rows
\usepackage{booktabs} % beautiful tables
\usepackage{graphicx, subfig} % figures
\usepackage{amsmath} % math
\usepackage{array, booktabs, tabularx}
\usepackage{animate} % animation of images
\usepackage{caption} % label
\captionsetup[figure]{font=scriptsize, justification=centering} % figure configuration
\setbeamerfont{footnote}{size=\tiny} % footnote configuration


% my colours
\definecolor{forestgreen}{HTML}{1d771d}
\definecolor{darkyellow}{HTML}{DAA520}
\definecolor{blue}{HTML}{007FFF}

% to make comment
\newcommand{\comment}[1]{}

% outline configuration
\setbeamerfont{section in toc}{size=\large}
\setbeamerfont{subsection in toc}{size=\small}
\setbeamertemplate{section in toc}[sections numbered]

% remove buttons of Madrid theme
\setbeamertemplate{navigation symbols}{}

% to remove caption from subfloat
\captionsetup[subfigure]{labelformat=empty}


%\newcommand{\mcx[1]}{\multicolumn{1}{>{\centering\arraybackslash}X}{\textbf{#1}}}	



% introduction parameters
\title[Laboratorio 0]{Control no lineal}
\subtitle{\textbf{Introducción a los módulos de control y Simulink}}

\author[J. Charaja and R. Terreros (UTEC)]{\textbf {Profesora} \\ Ruth Canahuire \\ \vspace{1em} \textbf{Asistentes de enseñanza} \\ Jhon Charaja y Ricardo Terreros}

\logo{\includegraphics[height=8em]{imgs/utec-logo.png}}

\date{\today}


\begin{document}
	\frame{\titlepage}
	\section{Introducción}
	% hablar sobre los módulos que vamos a usar	
	\begin{frame}
		\frametitle{Módulos del laboratorio L414}
		
		\begin{minipage}[t]{0.47\textwidth}
			\graphicspath{{imgs/qube/}}
			\centering
			\animategraphics[controls={play}, loop, width=\textwidth]{60}{}{10}{20} \\
			\vspace{1em}
			{\large \textbf{QUBE-Servo 2\footnotemark[1]}}
		\end{minipage}
		%\hspace{.08\textwidth}
		\begin{minipage}[t]{0.47\textwidth}
			\graphicspath{{imgs/inverted_pendulum/}}
			\centering
			\animategraphics[controls={play}, loop, width=\textwidth]{60}{}{10}{20} \\
			\vspace{1em}
			{\large \textbf{Inverted Pendulum\footnotemark[2]}}
		\end{minipage}
				
		\footnotetext[1]{Video recuperado de \url{https://www.quanser.com/products/qube-servo-2/} }
		\footnotetext[2]{Video recuperado de \url{https://www.quanser.com/products/qube-servo-2/} }
	\end{frame}
\end{document}